\documentclass[../../main.tex]{subfiles}
% \documentclass[a4paper,10pt]{ltjsarticle}
\begin{document}
\rhead{京都大学後期$1996$年$2$番}
\setcounter{equation}{0}
\begin{oframed}
  \subfile{problem}
\end{oframed}
\setlength{\columnseprule}{0.4pt}
\begin{multicols}{2}
  {\bf [解答]}
  表示の簡潔さのため
  \begin{align}
    A & =m^n+1 \\
    B & =n^m+1
  \end{align}
  とおく．ただし
  \begin{align}
    & m,n \in \mathbb{N} \\
    & m<n
  \end{align}
  とする．

  $n,m$がともに偶数とすると，$\mod 2$で考えると
  \begin{align}
    A & \equiv 1 \pmod{2}\\
    B & \equiv 1 \pmod{2}
  \end{align}
  となり，$A$も$B$も$10$の倍数にならず不適である．

  次に$n,m$のうち片方が奇数，もう片方が偶数の場合を考える．
  例えば$m$が偶数で$n$が奇数とすると，
  \begin{align}
    A \equiv 1 \pmod{2}
  \end{align}
  となり$A$が$10$の倍数にならずに不適である．逆の場合も対称性から同様である．

  そこで$n,m$は共に奇数であり，この時$A,B$は
  \begin{align}
    A & = m^n+1 = (m+1)(m^{n-1}-m^{n-2}+\dots+1) \\
    B & = n^m+1 = (n+1)(n^{m-1}-n^{m-2}+\dots+1)
  \end{align}
  のように分解できる．
  これらが共に$10$の倍数であるものの一つに，
  $m+1, n+1$が共に$10$の倍数であるような組がある．
  $m<n$から．$m=9, n=19$とすれば良い．  $\cdots$(答)

  \vspace{10pt}
  {\bf [解説]}
  最初から$\mod 10$で考えても良い．$A,B$が$10$の倍数である時，
  \begin{align}
    m^n \equiv -1 \\
    n^m \equiv -1
  \end{align}
  となる．一例として$(m,n)=(9,19)$とすると
  \begin{align}
    (-1)^{19} \equiv -1 \\
    (-1)^9 \equiv -1
  \end{align}
  となり条件を満たす．

  \newpage
\end{multicols}
\end{document}
